\noindent\textit{Poznámka: níže uvedené body jsou obecná doporučení (``general background knowledge'') a měly by být ověřeny lokálním pilotem před plošným nasazením.}\medskip
Cíle automatizovaného hodnocení kódu: rychlá verifikace funkčnosti, konzistentní kontrola chyb a okamžitá zpětná vazba studentům, přičemž učitelé se mohou soustředit na design a komplikovanější hodnocení.
Klíčové komponenty: testovací sada (unit/integrace, chybové stavy), statická analýza, metriky kvality a lidský dohled; návrh nasazení: jasná rubrika (co automaticky vs. ručně), validační anonymní vzorek, workflow odevzdání \(\to\) automatické testy \(\to\) učitelova kontrola a malý pilot.
Hodnocení systému: shoda s ručním hodnocením, flaky tests, pokrytí a míra nutného lidského zásahu; tipy: používat automatiku pro repetitivní úlohy, dokumentovat hranice, řešit fluktuaci testů a zvážit nástroje pro detekci plagiátorství; rychlý checklist: rubrika, testy/lintery, validační vzorek, komunikace se studenty, plán auditu.